%Theory
\section{Literature}
The Vehicle Routing Problem (VRP) belongs to the class of NP-hard problems and represents one of the most extensively studied topics in logistics and transportation. Its basic version, first proposed by Dantzig and Ramser in 1959, assumes the optimization of vehicle routes serving a set of customers while minimizing the total travel cost \cite{Dantzig1959}. Since then, many extensions have been developed, incorporating constraints such as vehicle capacity (CVRP), service times, different vehicle types, and time windows within which customers can be served — the so-called Vehicle Routing Problem with Time Windows (VRPTW) \cite{Toth2014, Golden2008}. The capacitated version with time windows (CVRPTW) better reflects the conditions of real-world transportation systems, where resource availability and time requirements play a crucial role.

In modern times, logistics systems operate under increasingly dynamic conditions — not only in terms of variable transportation times and fluctuating fleet sizes but also regarding demand and changing time windows within which deliveries and unloading must occur. For example, in local grocery stores with fresh products, goods must be delivered continuously and in appropriate quantities to ensure freshness and availability while avoiding spoilage. Another example is the construction industry, where material deliveries must be synchronized with the work schedule — in such cases, any mismatch may lead to delays, downtime, or extended working hours, resulting in costly losses. These situations require routing models capable of accounting for uncertainty and dynamic variability in data.

Traditional VRPTW and CVRPTW models assume deterministic data — fixed demand, travel times, and rigid time windows \cite{surveyBraekers, RCVRPTWreviewYernar}. In reality, however, transportation systems operate in environments where these parameters frequently change, making deterministic solutions less effective in practice. Consequently, there is a growing interest in robust and stochastic models that account for the variability of operational data.

On one hand, the literature on robust variants of the VRP has been developing intensively — reviews indicate numerous studies addressing random demand and stochastic travel times \cite{Oyola2018, Bomboi2022}. One of the works that implements a VRPTW model with uncertain data, where requests arrive randomly and time windows start according to the moment the demand appears, is the study by Pavone et al. \cite{Pavone2009}. However, in most such studies, time windows remain deterministic, even though in real-world systems they can shift dynamically.

On the other hand, in the field of robust optimization, research has been conducted to safeguard the routing process against data uncertainty. For example, Munari et al. \cite{Munari2018} proposed an RVRPTW model considering uncertain demand and travel times, along with a branch-price-and-cut algorithm. Zhang and Shen \cite{Zhang2021} developed a data-driven distributionally robust model for the VRPTW with time windows, minimizing the risk of delays under unknown travel time distributions. Similarly, Shen et al. \cite{Shen2022} investigated electric vehicle routing problems where uncertain parameters included demand and energy consumption dependent on load weight and time, demonstrating that uncertainty must also be considered for electric vehicles.

Despite the existing studies in the literature, significant research gaps remain. Many works treat random demand or uncertain travel times as stochastic variables, while only a small number of studies consider variable time windows. For example, Bomboi et al. \cite{Bomboi2022} showed that in the VRPTW model with random travel times, uncertainty strongly affects route feasibility. Errico et al. \cite{Errico2016} analyzed the VRPTW with stochastic service times but under deterministic time windows. On the other hand, Remy Spliet and A.F. Gabor \cite{Spliet2014} introduced the Time Window Assignment Vehicle Routing Problem (TWAVRP), in which time windows themselves are decision variables determined before demand is realized, thus taking an important step toward routing under uncertainty. Moreover, there is almost no research combining variable time windows and variable demand, which would extend the VRPTW to its capacitated form, CVRPTW. Another notable gap concerns scenario-based approaches, where multiple randomly perturbed instances are generated and solved independently, allowing the assessment of plan robustness. Thus, despite new studies emerging — for instance, Serrano et al. \cite{Serano2024} analyzed the VRPTW with travel times dependent on contextual features — a clear research gap still exists in the area of CVRPTW with variable time windows and stochastic demand.

The systematic review by Liu et al. \cite{Liu2023} indicates that approximately 86\% of studies on VRPTW employ approximate methods (heuristics and metaheuristics), with the vast majority being tested on deterministic data. This implies that, despite the widespread use of advanced search techniques, research on VRPTW under uncertainty still represents a minority within the field.

In summary, despite the notable progress achieved in the literature, a clear research gap can still be identified in the following areas:
\begin{itemize}
    \item Modeling random time window boundaries as variables drawn from probability distributions.
    \item Integrating the stochastic nature of time windows with random demand.
    \item Applying scenario-based optimization approaches and analyzing results in terms of robustness, stability, and risk (e.g., the occurrence of extreme cost scenarios).
    \item Conducting experiments on large benchmark instances that combine the above elements for future research.
\end{itemize}
This study proposes an approach that aims to fill the identified research gap. It allows for a better understanding of the problem and provides a foundation for further research, such as comparative studies or the development of new solution methods.
